\documentclass[12pt,a4paper]{article}
\usepackage[margin=25mm, headheight=15pt, headsep=10pt]{geometry}
\usepackage{fontspec}
\usepackage[table]{xcolor} % Note: table option is required for rowcolors
\usepackage{titlesec}
\usepackage{tocloft}
\usepackage{fancyhdr}
\usepackage{booktabs}
\usepackage{tcolorbox}
\tcbuselibrary{skins, breakable}
\usepackage[colorlinks=true, linkcolor=emerald, urlcolor=emerald, bookmarks=true]{hyperref}

\definecolor{emerald}{RGB}{0,102,51}
\definecolor{darkred}{RGB}{139,0,0}
\definecolor{lightgray}{RGB}{245,245,245}

\usepackage[nil]{babel}
\babelprovide[import, main]{bengali}
\babelprovide[import]{arabic}
\babelfont{rm}[Renderer=HarfBuzz,Scale=1.0]{Noto Serif Bengali}
\babelfont{sf}[Renderer=HarfBuzz]{Noto Sans Bengali}
\babelfont{tt}[Renderer=HarfBuzz]{Noto Sans Bengali}
\babelfont[arabic]{rm}[Renderer=HarfBuzz,Scale=1.1]{Noto Naskh Arabic}

% Section styling
\titleformat{\section}{\Large\bfseries\color{emerald}}{\thesection.}{0.5em}{}
\titleformat{\subsection}{\large\bfseries\color{emerald!85!black}}{\thesubsection.}{0.5em}{}
\titleformat{\subsubsection}{\normalsize\bfseries\color{emerald!70!black}}{\thesubsubsection.}{0.5em}{}
\titlespacing*{\section}{0pt}{18pt}{8pt}

% Fancy Header/Footer setup
\pagestyle{fancy}
\fancyhf{} % clear all header and footer fields
\fancyhead[L]{\color{emerald}\small\textbf{\nouppercase{\leftmark}}}
\fancyfoot[L]{\scriptsize\color{black!50}{\lat{AI}}-সংকলিত, আলেম-যাচাইকৃত নয়}
\fancyfoot[C]{\color{emerald!70!black}\thepage}
\renewcommand{\headrulewidth}{0.4pt}
\renewcommand{\headrule}{\hbox to\headwidth{\color{emerald}\leaders\hrule height \headrulewidth\hfill}}
\renewcommand{\footrulewidth}{0pt}

% Custom boxes
% 1. Ayat/Hadith Box
\newtcolorbox{ayatbox}[1][]{
    enhanced, breakable, 
    colback=emerald!5, 
    colframe=emerald, 
    boxrule=0.5pt, 
    arc=3pt, 
    left=10pt, right=10pt, top=10pt, bottom=10pt,
    #1
}

% 2. Quote/Translation Box
\newtcolorbox{quotebox}[1][]{
    enhanced, breakable,
    frame hidden,
    colback=lightgray,
    borderline west={3pt}{0pt}{emerald},
    left=15pt, right=10pt, top=10pt, bottom=10pt,
    sharp corners,
    #1
}

% 3. AI-generated notice box
\newtcolorbox{warnbox}[1][]{
    enhanced, breakable,
    colback=red!4,
    colframe=darkred,
    boxrule=0.8pt,
    arc=3pt,
    left=10pt, right=10pt, top=10pt, bottom=10pt,
    #1
}

% Commands
\newcommand{\arab}[1]{\foreignlanguage{arabic}{#1}}
\newcommand{\kib}[1]{\textcolor{emerald}{\textbf{#1}}}

% Latin (English) fallback: Noto Serif Bengali has NO Latin glyphs
\newfontfamily\latfont[Renderer=HarfBuzz]{Noto Serif}
\newcommand{\lat}[1]{{\latfont #1}}

\setlength{\parskip}{5pt}
\setlength{\parindent}{0pt}

% Fix for missing bullet in Noto Serif Bengali
\renewcommand{\labelitemi}{$\bullet$}



\begin{document}
\begin{center}{\Large\bfseries\color{emerald} \lat{14}-কুনূত-নফল-বিশেষ-দোয়া}\end{center}
\vspace{1em}
\section*{১৪ — কুনূত, নফল ও বিশেষ দোয়া}

\begin{warnbox}
 \textbf{গুরুত্বপূর্ণ নোটিশ:} এই কন্টেন্টটি \lat{AI} দিয়ে প্রস্তুত। কেবল সহীহ/হাসান হাদিস রাখা হয়েছে।
\end{warnbox}

\medskip\hrule\medskip

\subsection*{১. কুনূত (ওয়িতরে পড়া হয়)}

\subsubsection*{মূল আরবি}
\begin{center}\arab{اللَّهُمَّ اهْدِنِي فِيمَنْ هَدَيْتَ، وَعَافِنِي فِيمَنْ عَافَيْتَ، وَتَوَلَّنِي فِيمَنْ تَوَلَّيْتَ، وَبَارِكْ لِي فِيمَا أَعْطَيْتَ، وَقِنِي شَرَّ مَا قَضَيْتَ، فَإِنَّكَ تَقْضِي وَلَا يُقْضَىٰ عَلَيْكَ، وَإِنَّهُ لَا يَذِلُّ مَنْ وَالَيْتَ، وَلَا يَعِزُّ مَنْ عَادَيْتَ، تَبَارَكْتَ رَبَّنَا وَتَعَالَيْتَ}\end{center}

\subsubsection*{বাংলা উচ্চারণ}
\textbf{আল্লাহুম্মাহদিনী ফীমান হাদাইত, ওয়া'আফিনী ফীমান 'আফাইত, ওয়া তাওয়াল্লানী ফীমান তাওয়াল্লাইত, ওয়া বারিক লী ফীমা আ'তাইত, ওয়া কুনি শাররা মা কায়াইত, ফাইন্নাকা তাক্বদী ওয়ালা ইউক্বদা 'আলাইক, ওয়া ইন্নাহু লা যাইল্লু মান ওয়ালাইত, ওয়ালা ইয়া'যুয্যু মান 'আদাইত, তাবারাকতা রাব্বানা ওয়া তা'আলাইত}

\subsubsection*{শব্দ-শব্দ অর্থ}
\begin{table}[h]\centering\small
\begin{tabular}{lll}
আরবি & বাংলা & \lat{English} \\
\hline
\textbf{\arab{اهْدِنِ}ি} & আমাকে হেদায়াত দাও & \textit{\lat{guide} \lat{me}} \\
\textbf{\arab{فِيمَنْ} \arab{هَدَيْتَ}} & যাদের তুমি হেদায়াত দিয়েছ & \textit{\lat{among} \lat{those} \lat{You} \lat{guided}} \\
\textbf{\arab{وَعَافِنِي}} & আমাকে আরোগ্য দাও & \textit{\lat{grant} \lat{me} \lat{well-being}} \\
\textbf{\arab{وَتَوَلَّنِ}ি} & আমার তরফদারি কর & \textit{\lat{be} \lat{my} \lat{Guardian}} \\
\textbf{\arab{وَبَارِكْ} \arab{لِي}} & আমার জন্য বরকত দাও & \textit{\lat{bless} \lat{for} \lat{me}} \\
\textbf{\arab{وَقِنِي} \arab{شَرَّ}} & অনিষ্ট থেকে রক্ষা কর & \textit{\lat{protect} \lat{me} \lat{from} \lat{evil}} \\
\textbf{\arab{تَبَارَ}ক\arab{ْتَ}} & তুমি বরকতময় & \textit{\lat{Blessed} \lat{are} \lat{You}} \\
\textbf{\arab{تَعَالَيْتَ}} & তুমি মহান & \textit{\lat{Exalted} \lat{are} \lat{You}} \\
\hline
\end{tabular}
\end{table}

\subsubsection*{পূর্ণ বাংলা অনুবাদ}
\textbf{\lat{“}হে আল্লাহ! যাদের তুমি হেদায়াত দিয়েছ তাদের মধ্যে আমাকেও হেদায়াত দাও, যাদের তুমি আরোগ্য দিয়েছ তাদের মধ্যে আমাকেও আরোগ্য দাও, যাদের তুমি তরফদারি করেছ তাদের মধ্যে আমাকেও তরফদারি কর, যা দিয়েছ তাতে বরকত দাও, যা বিচার করেছ তার অনিষ্ট থেকে রক্ষা কর। নিশ্চয়ই তুমি বিচার করেন, তোমার বিপরীতে বিচার হয় না। তোমার অনুগৃহীত লজ্জিত হয় না, তোমার শত্রু সম্মানিত হয় না। তুমি বরকতময়, মহান।\lat{”}}

\medskip\hrule\medskip

\subsection*{২. ইস্তিখারা দোয়া (বিশেষ সিদ্ধান্তের আগে)}

\subsubsection*{মূল আরবি}
\begin{center}\arab{اللَّهُمَّ إِنِّي أَسْتَخِيرُكَ بِعِلْمِكَ وَأَسْتَقْدِرُكَ بِقُدْرَتِكَ وَأَسْأَلُكَ مِنْ فَضْلِكَ الْعَظِيمِ فَإِنَّكَ تَقْدِرُ وَلَا أَقْدِرُ وَتَعْلَمُ وَلَا أَعْلَمُ وَأَنْتَ عَلَّامُ الْغُيُوبِ}\end{center}

\subsubsection*{পূর্ণ বাংলা অনুবাদ}
\textbf{\lat{“}হে আল্লাহ! আমি তোমার জ্ঞান দিয়ে কল্যাণ চাই, তোমার ক্ষমতা দিয়ে শক্তি চাই এবং তোমার মহান অনুগ্রহ থেকে চাই — কারণ তুমি ক্ষমতাবান, আমি নই; তুমি জানেন, আমি জানি না; এবং তুমি অদৃশ্য বিষয়ের জ্ঞাতা।\lat{”}}

\medskip\hrule\medskip

\subsection*{৩. তাসবিহ ফাতিমা — প্রয়োগ}

\begin{table}[h]\centering\small
\begin{tabular}{ll}
ধাপ & কাজ \\
\hline
১ & নামাজের পর বসে ডান হাতের আঙুল দিয়ে গোনুন \\
২ & \textbf{সুবহানাল্লাহ} ৩৩ বার \\
৩ & \textbf{আলহামদুলিল্লাহ} ৩৩ বার \\
৪ & \textbf{আল্লাহু আকবার} ৩৩ বার \\
৫ & \textbf{লা ইলাহা ইল্লাল্লাহ...} ১ বার — মোট ১০০ \\
\hline
\end{tabular}
\end{table}

\medskip\hrule\medskip

\subsection*{৪. সংশ্লিষ্ট হাদিস}

\begin{table}[h]\centering\small
\begin{tabular}{ll}
হাদিস & গ্রন্থ \\
\hline
\textbf{\lat{“}নামাজের পর 'সুবহানাল্লাহ' ৩৩ বার, 'আলহামদুলিল্লাহ' ৩৩ বার, 'আল্লাহু আকবার' ৩৩ বার বলো।\lat{”}} & সহীহ মুসলিম ৫৯১ \\
\textbf{\lat{“}কুনূত ওয়িতরের শেষ রাকাতে পড়া হয়।\lat{”}} & সুনানে আবু দাউদ ১৪২৫ \\
\textbf{\lat{“}ইস্তিখারা দোয়া জানাশনায় শিখিয়ে দেওয়া হয়েছে।\lat{”}} & সহীহ বুখারি ১১৬২ \\
\hline
\end{tabular}
\end{table}

\medskip\hrule\medskip

\textit{শেষ আপডেট: আগস্ট ২০২৬}


\end{document}
